\section{Taxonomy of Excusable vs. Non-Excusable Rust Behavior}
\label{sec:bad_rust}

Rust provides developers with highly robust memory safety and concurrency guarantees. However, it also includes tools such as \texttt{unsafe}, \texttt{unwrap}, macros, and \texttt{Arc<Mutex<T>>} that are required for systems programming but can easily be abused. In the era of ``vibe coding''---where developers rely on Large Language Models (LLMs) to iteratively patch compiler errors---the temptation to brute-force a green build often overrides formal correctness. This section establishes a taxonomy of ``never justified'' behavior in production Rust code, distinguishing between subjective stylistic choices and objective correctness, safety, or security failures \cite{rust_reference_ub, rustonomicon}.

\subsection{The ``Never Justified'' Baseline}

Bad Rust behavior is never justified when it satisfies one of the following conditions:
\begin{enumerate}
    \item \textbf{Introduces Undefined Behavior (UB):} Violates aliasing rules, dereferences dangling pointers, or produces invalid values.
    \item \textbf{Lies to the Type System:} Exposes a safe API that secretly requires unsafe preconditions to avoid UB.
    \item \textbf{Unproven \texttt{unsafe} Code:} Uses \texttt{unsafe} without a formally written, local, reviewable \texttt{SAFETY:} argument \cite{std_dev_guide_safety}.
    \item \textbf{Misuses Panics:} Turns expected failures into process panics rather than using \texttt{Result}.
    \item \textbf{Blind Trust in Generation:} Ships AI/generated code or opaque dependencies that the author does not fully comprehend \cite{arxiv_llm_crypto}.
\end{enumerate}

\begin{table}[htbp]
    \centering
    \caption{Top Ranked Non-Excusable Rust Anti-Patterns}
    \label{tab:bad_rust_top_issues}
    \renewcommand{\arraystretch}{1.3}
    \begin{tabular}{|c|l|p{8.5cm}|}
        \hline
        \textbf{Rank} & \textbf{Category} & \textbf{Anti-Pattern \& Impact} \\
        \hline
        1 & Memory Safety & \textbf{Unsound Safe APIs \& UB:} Exposing APIs as safe when a caller can trigger Undefined Behavior compromises the ecosystem trust model. \\
        \hline
        2 & \texttt{unsafe} Abuse & \textbf{Undocumented \texttt{unsafe}:} Using \texttt{unsafe} without a \texttt{SAFETY:} comment that proves invariants reduces Rust to C with worse tooling. \\
        \hline
        3 & Concurrency & \textbf{Clone/\texttt{Arc<Mutex>} Defaults:} Blindly cloning or wrapping types to bypass the borrow checker obscures ownership and introduces deadlocks. \\
        \hline
        4 & Error Handling & \textbf{Unjustified \texttt{unwrap()}:} Crashing on expected, external, or attacker-controlled input turns routine errors into DoS vulnerabilities. \\
        \hline
        5 & FFI Boundaries & \textbf{Unsound C Interop:} Passing raw Rust references across the FFI boundary or trusting C data without validation instantly breaks memory safety. \\
        \hline
        6 & Async Rust & \textbf{Blocking the Executor:} Holding blocking locks or doing heavy synchronous I/O across \texttt{.await} points starves executors. \\
        \hline
        7 & Vibe Coding & \textbf{``Make it Compile'' Brute Force:} Accepting AI-generated code by adding \texttt{unwrap} or \texttt{clone} until it compiles, without semantic understanding. \\
        \hline
        8 & Supply Chain & \textbf{Blind Dependencies:} Adding unreviewed crates or ignoring RustSec advisories introduces immediate supply chain risks. \\
        \hline
        9 & API Design & \textbf{``Stringly'' Typed Domains:} Using primitive types for validated states makes illegal states representable and pervasive. \\
        \hline
        10 & Security & \textbf{In-Band Secrets:} Rolling custom cryptography, logging secrets, or using weak RNGs bypasses language-level safety. \\
        \hline
    \end{tabular}
\end{table}

\subsection{Detailed Taxonomy of Bad Behavior}

\paragraph{1. Memory Safety and \texttt{unsafe} Abuse.}
The most severe violations involve creating data races, bypassing aliasing rules, and reading uninitialized memory \cite{rust_reference_ub}. Using \texttt{unsafe} is not inherently bad, but doing so to bypass the borrow checker or without rigorous \texttt{SAFETY:} proofs is inexcusable. This includes exposing safe APIs that can trigger UB, or hiding \texttt{unsafe} blocks inside macros.

\paragraph{2. FFI, ABI, and Interoperability Sins.}
The Foreign Function Interface (FFI) is where Rust's guarantees vanish. Inexcusable behaviors include passing Rust references to C that outlive the call, using incorrect ABIs, guessing C struct layouts, and allowing unwinding across FFI boundaries that do not support it.

\paragraph{3. Error Handling and Panic Misuse.}
While \texttt{unwrap()} and \texttt{expect()} are acceptable in tests and prototypes, their use on user input, network data, or in production libraries is unjustifiable \cite{rust_api_guidelines}. Expected failures must be handled via \texttt{Result}, and turning them into panics constitutes a denial-of-service risk.

\paragraph{4. Async and Concurrency Anti-Patterns.}
Rust's async ecosystem requires specific care. Holding a \texttt{std::sync::MutexGuard} across an \texttt{.await} point, blocking the runtime with heavy CPU/IO workloads, and treating \texttt{tokio::spawn} as fire-and-forget without a shutdown plan are major violations \cite{tokio_docs}.

\paragraph{5. Type System and API Design.}
Rust's type system is designed to make illegal states unrepresentable. Failing to utilize this---by representing validated domain concepts as raw \texttt{String}s or \texttt{bool}s, exposing internal invariants via public fields, or implementing traits like \texttt{Deref} or \texttt{Eq} inconsistently---undermines the language's core strengths.

\paragraph{6. Supply Chain and Security Negligence.}
Accepting untrusted data without validation (e.g., deserializing into trusted domain types directly), rolling custom cryptography, or adding unreviewed crates heavily compromises system security. In the vibe-coding workflow, relying on LLMs for dependency selection without auditing for typosquatting or RustSec advisories is particularly hazardous \cite{rustsec}.

\paragraph{7. The Vibe-Coding Trap.}
The most pervasive modern anti-pattern is using the compiler as an interactive patch oracle. Developers or AI agents repeatedly apply \texttt{clone}, \texttt{Rc}, \texttt{RefCell}, \texttt{Box::leak}, or \texttt{unsafe} to silence borrow checker errors, resulting in code that compiles but fundamentally misrepresents the intended ownership and lifecycle model. This ``it compiles, ship it'' mentality replaces provable correctness with an illusion of safety.
